%
% Capítulo 1
%
\chapter{Introduction} \label{cap:introduction}

This chapter aims to present our project, describe all its functionalities, and explain the organization of our report.


%
% Secção 1.1
%
\section{Motivation} \label{sec11}

The objective of this project is to develop ClimbBeta, a full-stack platform dedicated to the management, logging, and sharing of indoor climbing and bouldering activities. The application targets two main audiences: practitioners (climbers) and sports facility managers (climbing gyms). Currently, the climbing community focuses heavily on technical progression and the memorization and sharing of route solutions (known as betas). However, logging these ascents and visually capturing problems are done sparsely across generic social networks. Simultaneously, gyms lack centralized tools to showcase their route catalogs and obtain metrics on their community's performance and preferences. This project aims to bridge this gap by offering a hybrid logbook. On one hand, it allows gym managers to map their walls and routes. On the other hand, it enables climbers to log their ascents either automatically (linked to a gym) or manually (outdoor rock climbing), accompanied by a multimedia visual diary and statistics sharing. The platform not only saves the climber time when logging data but also centralizes the social and statistical aspects into a single ecosystem.


%
% Secção 1.2
%
\section{System Requirements} \label{sec12}

In this subsection, we will address the functional and non-functional requirements of our
project.

%
% Secção 1.2.1
%
\subsection{Functional Requirements} \label{sec121}
\begin{itemize}
\item \textbf{Profile and User Type Management:} The system must support two account types: Climber (Regular User) and Gym Owner. Users must be able to register, authenticate, and manage their profiles.
\item \textbf{Gym Management (Profile: Gym Owner):} Coaches should be able to add and remove their athletes. Adding an athlete will only require their name and date of birth.
\begin{itemize}
\item \textbf Owners can register their facility.
\item \textbf Map the space into different walls.
\item \textbf On each wall, the Owner can add and remove Routes/Boulders.
\item \textbf Set base parameters: hold color, difficulty grade, creation date, and a route sketch/photograph.
\end{itemize}
\item \textbf{Ascent Logging (Hybrid Logbook - Profile: Climber):}
\begin{itemize}
\item \textbf{Integrated Logging (Gyms):} The climber can browse a registered gym's catalog, select a specific boulder, and the form will be auto filled with the route's official difficulty and details.
\item \textbf{Free Logging (Outdoor/Unlisted):} The system will allow manual and independent logging of ascents (e.g., outdoor bouldering or non-partner gyms), allowing the climber to define the location, sector name, and difficulty grade.
\item \textbf In both modes, the climber logs their performance (number of attempts, ascentstyle like Flash or Top, technical notes) and can attach photographs to document the Beta or location.
\end{itemize}
\item \textbf{Social Component and Favorites:} Functionality to search, follow, and add other climbers to a "Favorite Climbers" list (friends). Ability to save specific routes to a "Favorite Boulders / Projects" list to complete in the future.
\item \textbf{Statistics and Leaderboards:} Each boulder registered in a gym will have a dedicated page with global statistics and a community history. This page will display a ranking ordered by climbers' performance (e.g., who completed it Flash style, or an ordered list by the fewest attempts).
%
% Secção 1.2.2
%
\subsection{Non-Functional Requirements} \label{sec122}
\begin{itemize}
\item \textbf{Cloud Storage (Non-Functional):} Multimedia content (images) should not overload the relational database and must be stored in an external Cloud infrastructure (e.g., Firebase Storage or AWS S3), storing only the references (URIs) in the central system.
\end{itemize}

%
% Secção 1.2.3
%
\subsection{Optional Features} \label{sec123}
\begin{itemize}
\item \textbf{Video Support (Optional):} In an advanced phase, expanding the "Visual Diary" to support the upload and playback of short video demonstrations of the resolutions (Betas).
\item \textbf{Offline-First Mode (Optional):} Local support in the mobile application to allow browsing and partial submission of ascents even without internet access, mitigating the poor connectivity typical of indoor environments and outdoor climbing areas, synchronizing data later.
\end{itemize}


%
% Secção 1.3
%
\section{Report Organization} \label{sec13}
\begin{itemize}
The project report is organized into 5 chapters. The first corresponds to the Chapter 2,
Problem, where it is demonstrated how coaches record data in Excel sheets and some of the
terms used are explained. The second corresponds to the Chapter 3, Architecture, where the
data model and the technologies used are presented. The third corresponds to the Chapter 4,
Backend, where the organization and functionality of each layer of the API, HTTP, Services,
Repository, and Domain are explained, the authentication process, how the errors are handled
on the server-side and the communication with the database. The fourth corresponds to
the Chapter 5, Frontend, which describes the organization, communication with the API,
how errors are handled, and relative aspects to the web application and mobile application,
such as the organization, navigation and authentication process. Finally, in the Chapter 6,
Conclusion, the conclusions of this project are presented.
\end{itemize}

